\scanpage{187}
soit propre sur $X$. \deleted{Alors} Soit $V=E-g(X)$, où $g$ est la section unité, alors
\[
D\cap V\sim0\quad\text{dans }V.
\]

\diagram{\begin{tikzcd}[ampersand replacement=\&]
E\arrow[d,"f"']\& E'\arrow[l,"h'"']\arrow[d,"f'"]\\
X\& X'\arrow[l,"h"]\arrow[u,bend right=36,"q"']
\end{tikzcd}}
On peut supposer que $D$ est un diviseur irréductible, défini par un sous-préschéma fermé $X'$ de $E$. \deleted{\unsure{Faisons}}

Supposons d'abord $X'$ \uline{régulier}, faisons le changement de base $h:X'\longrightarrow X$, alors $E'$ sur $X'$ est muni d'une section naturelle $q$ telle que $h'q=$ immersion can. de $X'$ dans $X$. On a donc \deleted{dans Pic}\inserted{(raisonnant sur Pic\ldots]}
% slanted
\margin{N.B. Cela a un sens car $h'$ et $q$ sont finis\ldots}
\[
D=(h'_{*}q)_*(1_{X'})=h'_*(q_*{\unsure{q}}.1_{X'})
\]
Il est bien connu d'ailleurs que
\[
q_*(1_{X'})=f'{}^*(\xi_{E'})
\]
où $\xi_{E'}=\operatorname{Pic}(X')$ est la classe can. de $E'/X'$. D'ailleurs $\xi_{E'}=h^*(\xi_E)$, d'où
\[
\begin{aligned}
q_*(1_{X'})&=\deleted{\unsure{h}}f'{}^*h^*(\xi_E)=h'{}^*f^*(\xi_E)\qquad\text{et}\\
\unsure{{}'D|D}&=h'_*h'{}^*(f^*(\xi_E))=f^*(\xi_E)\,h'_*(1_{E'})
\end{aligned}
\]
% archive sheet number: 186
